%Ne pas numéroter cette partie
\part*{Annexes}
%Rajouter la ligne "Annexes" dans le sommaire
\addcontentsline{toc}{part}{Annexes}


\chapter*{Annexe 1}
\addcontentsline{toc}{chapter}{Annexe 1}

%changer le format des sections, subsections pour apparaittre sans le num de chapitre
\makeatletter
\renewcommand{\thesection}{\@arabic\c@section}
\makeatother

%recommencer la numérotation des section à "1"
\setcounter{section}{0}
Intro
\section{Partie 1}
Bla
\subsection{Sous-partie 1}
Bla
\subsection{Sous-partie 2}
Bla
\section{Partie 2}
Bla
\subsection{Sous-partie 1}
Bla

\chapter*{Annexe 2}
\addcontentsline{toc}{chapter}{Annexe 2}

%recommencer la numérotation des section à "1"
\setcounter{section}{0}
\begin{itemize}
\item item1;
\item item2;
\item item3;
\item item4.
\end{itemize}

Bla

\begin{center}
\textsc{Attention !}

\textit{Texte d'avertissement}
\end{center}

Bla

\newpage

\section{Partie 3}
Bla
% \begin{figure}[!ht]
% \begin{center}
% \includegraphics[height=8cm]{presentation/schema}
% \end{center}
% \caption[schema]{Presentation schema}
% \end{figure}

\paragraph*{Paragraphe 1}
~\\
\hskip7mm
Bla
\paragraph*{Paragraphe 2}
~\\
\hskip7mm
Bla

%tableau centré à taille variable qui s'ajuste automatiquement suivant la longueur du contenu
\begin{figure}[!h]
    \begin{center}
    \begin{tabular}{|l|l|l|l|l|}
      \hline
      Solution & Critère 1 & Critère 2 & Critère 3 & Critère 4\\
      \hline
      Solution 1(cf. ref. \cite{cite0}) & Oui & Oui & Oui & Oui \\
      Solution 2(cf. ref. \cite{cite1}) & Oui & Oui & Oui & Non \\
      Solution 3(cf. ref. \cite{cite2}) & Oui (sauf telle chose) & Non & Non & Oui\\
      Solution 4(cf. ref. \cite{cite3}) & Oui& Non & Oui & Non\\
      Solution 5(cf. ref. \cite{cite4}) & Oui (uniquement ceux-ci) & Non & Oui & Non\\
      \hline
    \end{tabular}
    \end{center}
    \caption{Tableau récapitulatif des solutions}
\end{figure}


%galerie d'image
% \begin{figure}[htp]
%     \centering
%     \subfloat[Première image]{\label{fig:première}\includegraphics[scale=0.8]{resultats/gallerie}}
%     ~ %espace entre deux images sur une même ligne
%     \subfloat[Deuxième image]{\label{fig:deuxième}\includegraphics[scale=0.8]{resultats/gallerie}}
%     ~
%     \subfloat[Troisième image]{\label{fig:troisième}\includegraphics[scale=0.8]{resultats/gallerie}}
%     ~\\ %saute une ligne dans la galerie d'image
%     \subfloat[Quatrième image]{\label{fig:quatrième}\includegraphics[scale=0.8]{resultats/gallerie}}
%     ~
%     \subfloat[Cinquième image]{\label{fig:cinquième}\includegraphics[scale=0.8]{resultats/gallerie}}
%     \caption{Différents screenshots quelque chose, en gallerie}
%     \label{fig:gallerie1}
% \end{figure}